% ===================================================================
%  नकसा नं. ७ — जाड़ में गाँव। बसंत में देहाती घर।
%
%  Traduction, faite en 2026, en bhojpouri d'aujourd'hui, sur l'ido
%  de texto/io. Regles de la colonne : voir 10-nakasa-01.tex. DEUX
%  scenes, deux numerotations qui repartent de 1 (t07-c1-*,
%  t07-c2-*). 53 blocs, 121 renvois, vingt-huit paires a retourner.
%  UNE note d'auteur, posee en tete. Le bloc t07-c1-07-1 tient DEUX
%  alineas sous une seule cle : ne pas en ouvrir une seconde.
%
%  LA SECONDE SCENE EST LA PREMIERE DU LIVRET QUI N'AIT DEMANDE AUCUN
%  REFUS. Six tableaux durant, cette colonne a du ecarter le mot le
%  plus juste de la langue parce que la chose n'etait pas la meme :
%  चइत pour les mois, राम-राम pour le salut, पतरा pour le calendrier,
%  कबड्डी pour le jeu de barres, रोटी pour le pain, खटिया pour le
%  lit. La ferme de l'alinea 2 a l'alinea 16 n'a rien demande de tel.
%  हर la charrue, नार le sillon, बीया la semence, इनार le puits, नाँद
%  l'auge, हेंगा la herse, खरिहान la grange, गोहाल l'etable, मथनी la
%  baratte, गोबर le fumier, गड़ेरिया le berger, हरवाह le laboureur :
%  tout se dit, et se dit du premier coup. UNE COUR DE FERME
%  FRANCAISE DE 1926 ET UN VILLAGE BHOJPOURI DE 2026 PARTAGENT ASSEZ
%  D'OBJETS POUR QUE LA TRADUCTION N'AIT PLUS RIEN A ARBITRER — et
%  c'est la premiere fois depuis le tableau 1.
%
%  MAIS LA FAUX N'A PAS DE NOM, ET C'EST LA COLONNE GUJARATIE QUI
%  L'AVAIT TROUVE LA PREMIERE. Elle avait releve que la faux n'est
%  jamais venue en Inde tandis que la faucille n'en est jamais
%  partie. Ici les « falchisti » (11) fauchent le pre avec leurs
%  « falchili » — et le bhojpouri n'a que हँसुआ, la faucille, qu'on
%  tient d'une main en s'accroupissant. On ecrit donc « बड़का हँसुआ »,
%  la grande faucille, faute de mieux. DEUX LANGUES INDIENNES SEPAREES
%  DE MILLE CINQ CENTS KILOMETRES, LE MEME TROU AU MEME ENDROIT : ce
%  n'est plus une particularite du gujarati, c'est une frontiere
%  d'outil.
%
%  ET LE MARRON (42) S'APPELLE « साहबलूत », LE CHENE-ROI. La
%  chataigne ne pousse pas dans la plaine, et la langue n'a pas
%  invente de mot : elle en a pris un au persan — شاه‌بلوط —, qui
%  nomme l'arbre par comparaison avec le chene. La colonne gujaratie
%  avait REFUSE la chataigne, faute de mot ; celle-ci en a un, et il
%  dit lui-meme qu'il vient d'ailleurs. Deux manieres de n'avoir pas
%  l'arbre.
%
%  Enfin le cochon (77, 78) s'ecrit सूअर, sans note et sans detour.
%  La colonne gujaratie avait rencontre le meme animal a son
%  tableau 7 — le meme numero — et avait du expliquer un interdit.
%  Ici la loi du releve joue sans commentaire : la langue a le mot,
%  c'est le faire qui se discute, pas le nommer.
% ===================================================================

%%K t07-noto-1 noto
\VUnotes{143.2mm}{%
(*) खाना के बिस्तार खातिर नकसा नं. ६ आ १३ देखीं।%
}

%%K t07-apar-1 apar
\VUsaut{4.0mm}
\VUcentre{13.2pt}{90}{दुसरका शृंखला}
\VUsaut{3.0mm}
\VUfilet{16mm}
\VUsaut{5.6mm}
\VUtitre{15.0pt}{60}{नकसा नं. ७}
\VUsaut{3.0mm}

%%K t07-tit-1 sub
\VUcentre{12.0pt}{0}{\textit{पहिलका दृस्य।}}
\VUsaut{3.0mm}

%%K t07-tit-2 sub
\VUpk{10.2pt}{40}{जाड़ में गाँव।}
\VUsaut{4.0mm}

%%K t07-c1-01-1 p
१. --- नया साल के छुट्टी में हम आपन बाबूजी के साथे नोवुरबो गइनी, एगो
छोट गाँव, जहाँ ओकरा कुछ बेपार के काम निपटावे के रहे।

%%K t07-c1-02-1 p
२. --- हमनी के ऊ \VUgras{डाक-गाड़ी} \textsuperscript{(11)} से गइनी जा जवन
सहर आ ओह कस्बा के बीच चलेला ; बाकिर बहुते ठंढा रहे, एहीसे ओह कम आराम
वाला गाड़ी के जातरा बहुते नीक ना रहल।

%%K t07-c1-03-1 p
३. --- एहीसे हमनी के सचहूँ खुसी भइल जब हमनी के \VUgras{गाड़ीवान} के
आपन गाड़ी चउक पर, \textit{उजर भालू} के \VUgras{सराय} \textsuperscript{(2)}
के आगे रोकत देखनी जा। \VUgras{सरायवाला} \textsuperscript{(4)} आपन घर के
देहरी पर, चुपचाप आपन \VUgras{पाइप} पियत, हमनी के अगोरत रहलें।

%%K t07-c1-03-2 p
ऊ हमनी के भीतर बोलवलें आ हम बिना दोसरा बेर कहवले ओह लकड़ी के आग के
सामने जा बइठनी जवन चूल्हा में चटकत रहे। तब हम ओह तेज \VUgras{उत्तर के
हवा} के खूब हँसनी जवन पुरनका \VUgras{तखती} \textsuperscript{(3)} के
चरमरावत रहे आ \VUgras{धुँआ} \textsuperscript{(1)} के, जवन अँगीठी से निकलत
रहे, करिया बवंडर बना के उड़ावत रहे।

%%K t07-c1-04-1 p
४. --- खाना के बेरा रहे। बिना जादे अगोरले, हमनी के ऊ सादा बाकिर
सवाद वाला खाना नीक से खइनी जा जवन हमनी के परोसल गइल (*)।

%%K t07-tit-3 sub
\VUpk{10.2pt}{40}{कुछ भेंट।}
\VUsaut{3.0mm}

%%K t07-c1-05-1 p
५. --- खाना के बाद हम आपन बाबूजी के साथे गइनी, जे बीमा के काम करेलें,
ओकर ग्राहक लोग लगे। पहिले हमनी के \VUgras{लोहार} \textsuperscript{(5)} लगे
गइनी जा, जे सराय के लगे रहेलें। ऊ घोड़ा के नाल ठोके में लागल रहलें। हम
ओकर संगी के मजबूती देख के दंग रह गइनी, जे घोड़ा के \VUgras{खुर} आपन
चमड़ा के फेंटा पर मजबूती से थामले रहलें, जबकि लोहार जोर-जोर के हथौड़ी
मार के \VUgras{नाल} \textsuperscript{(6)} बइठावत रहलें।

%%K t07-c1-06-1 p
६. --- ओकरा बाद हमनी के गाँव के \VUgras{मोचिया} \textsuperscript{(7)},
बुढ़उ याकोबुस, लगे गइनी जा। ओकर तखती पर सुंदर \VUgras{बूट} बनल रहे,
तबो ऊ फाटल जूता के जोड़ी में नया तला लगावे भर से संतोस कर लेलें।

%%K t07-c1-06-2 p
बुढ़उ हँसत हमनी के कहलें : « सहर में हम " जूता बनावे वाला " रहनी आ
हमरा लगे ग्राहक ना रहे। इहाँ हम " जूता सीये वाला " बानी आ काम के कवनो
कमी नइखे। »

%%K t07-c1-07-1 p
७. --- ऊ हमनी के आपन पड़ोसी \VUgras{हजाम} \textsuperscript{(8)} के बारे
में बतवलें, जेकरा ग्राहक लोग खुस नइखे। ओकर \VUgras{उस्तरा} हरमेसा
खँड़ल रहेला आ ओकर \VUgras{कैंची} कबो नीक से ना काटे।

बाकिर काहेकि ऊ गाँव के एकलौता हजाम हउवें, लोग के मजबूरी बा कि ओकरे से
हजामत करावे आ बाल कटावे। ऊपर से ऊ कंजूस आ रूखा मनई हउवें।

%%K t07-c1-08-1 p
८. --- गनीमत बा कि बाकी \VUgras{पड़ोसी} लोग ओकरा नियर नइखे। जइसे
\VUgras{गाड़ी बनावे वाला} \textsuperscript{(9)} भल मनई हउवें, मेहनती आ
मददगार। ओकरा से गाड़ी बनवावल सुख के बात बा। ओकर काम हरमेसा पूरा रहेला।

%%K t07-c1-09-1 p
९. --- बुढ़उ याकोबुस \VUgras{जीन बनावे वाला} \textsuperscript{(10)} के
बारे में भी सिकायत ना कइलें, जेकरा ऊ कबो-कबो \VUgras{साज} सीये में मदद
करेलें जब काम के जल्दी रहेला।

%%K t07-c1-09-2 p
हमार बाबूजी ओकरा से ओकर घर-परिवार के हाल पूछलें : « सब नीक बा। हमार
नतिनी \VUgras{इसकूल} \textsuperscript{(12)} में बाड़ी। ओकर
\VUgras{मास्टरनी} \textsuperscript{(13)} ओकरा से बहुते खुस बाड़ी, ऊ नीक
\VUgras{बिद्यार्थी} \textsuperscript{(14)} बाड़ी। ऊ हमार निकम्मा बेटा
नियर नइखी ; ओकरा खाली खेले के सूझेला। \VUgras{मास्टर साहेब}
\textsuperscript{(15)} ओकरा के कुछ भी ना सिखा पावेलें। »

%%K t07-tit-4 sub
\VUsaut{3.0mm}
\VUpk{10.2pt}{40}{गाँव के गपसप।}
\VUsaut{3.0mm}

%%K t07-c1-10-1 p
१०. --- ओकरा बाद बुढ़उ हमनी के गाँव के खबर सुनवलें। नया इसकूल बने वाला
बा, काहेकि जवन \VUgras{पंचायत घर} में बा ऊ बहुते छोट पड़ गइल। ई काम
आसान बा, काहेकि \VUgras{पनचक्की वाला} के बगइचा के एगो हिस्सा खरीदल
बहुते बा। ओह बेचारा खातिर ई बहुते फायदा के बात होई। ठंढ के चलते
\VUgras{पनचक्की} \textsuperscript{(16)} के \VUgras{पाट} आउर गोहूँ ना पीस
पावत बा, काहेकि \VUgras{चाक} \textsuperscript{(17)} के \VUgras{बरफ}
\textsuperscript{(25)} जमा के रोक देले बा, जवन \VUgras{नाला}
\textsuperscript{(23)} के ह।

%%K t07-c1-11-1 p
११. --- « बनावट के बात बा त, रउआ हमनी के नया \VUgras{गिरजाघर}
\textsuperscript{(18)} देखनी ? ऊ आपन बरफ के चादर के नीचे बहुते सुंदर
लागत बा। \VUgras{कउआ} \textsuperscript{(19)}, जे आपन पुरनका घंटाघर के आउर
ना पहिचानत बाड़ें, उदास होके \VUgras{कबरिस्तान} \textsuperscript{(20)} के
बड़का गाछ पर बइठल बाड़ें। »

%%K t07-c1-12-1 p
१२. --- कबरिस्तान के ई बात से मोचिया के ऊ लोग इयाद आ गइल जेकरा के
हमार बाबूजी जानत रहलें आ जे पिछला साल गुजर गइल। « केतना \VUgras{कबर}
\textsuperscript{(21)}, हमार भल साहेब, \VUgras{सरो} \textsuperscript{(22)}
के नीचे बाकी सब में जुड़ गइल ! मउअत हमनी के पुरनका नोटरी के काट लेलस।
पूरा गाँव ओकर \VUgras{अंतिम जातरा} में सामिल भइल। »

%%K t07-c1-13-1 p
१३. --- « ई रहलें ओकर उत्तराधिकारी, जे नाला पर फिसलत बाड़ें। ऊ अभिए
हमरा लगे आपन \VUgras{बरफ के जूता} \textsuperscript{(26)} के फीता बनवावे
आइल रहलें, जवन टूट गइल रहे। »

%%K t07-c1-14-1 p
१४. --- « ई रहलें नाला के किनारे नया \VUgras{खेत के चौकीदार}
\textsuperscript{(27)}। ऊ पुरनका सिपाही हउवें जे गाँव के सैतान लइकन के
डरा देलें। ऊ लोग ओकर \VUgras{गेटर} \textsuperscript{(29)} आ ओकर
\VUgras{कैप} \textsuperscript{(28)} देखते भाग जाला। »

%%K t07-c1-15-1 p
१५. --- « हा ! ई रहलें बुढ़उ योसेफ जे नानबाई खातिर \VUgras{लकड़ी के
गट्ठर के गाड़ी} \textsuperscript{(30)} ले जात बाड़ें। --- सही बा, हमार
बाबूजी कहलें, हम ओकर \VUgras{खच्चर} \textsuperscript{(31)} के जोड़ी के
पहिचानत बानी। हमरा ओकरा से बतियावे के बा, हम एह मउका के फायदा उठाइब।
फेरु भेंट होई, याकोबुस बाबा। »

%%K t07-c1-16-1 p
१६. --- हमनी के बाहर निकलते, गाँव के लइकन इसकूल से निकल के दउड़त
\VUgras{फिसलनी} \textsuperscript{(33)} ओर लपकल, ई खतरा उठा के कि
\VUgras{गिरल} \textsuperscript{(34)} में हाथ-गोड़ टूट जाई। ओहमें से एगो
गाड़ी बनावे वाला के इहाँ से आपन \VUgras{बरफ-गाड़ी} \textsuperscript{(32)}
लेलस, ओहमें आपन एगो संगी के बइठवलस आ दउड़त चल गइल।

%%K t07-c1-17-1 p
१७. --- दोसर लोग \VUgras{बरफ के ढेर} \textsuperscript{(37)} ओर लपकल,
जवन \VUgras{पुल} \textsuperscript{(40)} लगे रहे, आ ओह \VUgras{बरफ के
मूरत} \textsuperscript{(39)} पर गोला बरसावे लागल जवन ऊ लोग कल्ह बनवले रहे
आ जेकरा माथ पर ऊ लोग टोपी, पाइप आ कान्ह में लटकल इसकूल के झोरा धर देले
रहे।

%%K t07-c1-18-1 p
१८. --- ओह लोग के जमा देवे वाला हवा आ ठंढ के परवाह कम बा। जादेतर लोग
लगे त \VUgras{गुलूबंद} \textsuperscript{(35)} भी नइखे, आ \VUgras{बरफ के
गोला} \textsuperscript{(38)} के लड़ाई के बाद फेरु गरम होखे खातिर ओह लोग
के मुट्ठी भर बहुते गरम \VUgras{साहबलूत} \textsuperscript{(42)} बहुते बा,
जवन बुढ़िया \VUgras{बेचनिहारिन} याकोबिना से खरीदल जाला, जे आपन बहुते
पातर \VUgras{साल} \textsuperscript{(43)} के नीचे ठंढ से काँपत बाड़ी।

%%K t07-tit-5 sub
\VUsaut{3.0mm}
\VUcentre{12.0pt}{0}{\textit{दुसरका दृस्य।}}
\VUsaut{3.0mm}

%%K t07-tit-6 sub
\VUpk{10.2pt}{40}{बसंत में देहाती घर।}
\VUsaut{4.0mm}

%%K t07-c2-01-1 p
१. --- जाड़ के सब सुख के बादो, हमनी के गहिर खुसी से \VUgras{बसंत} के
लवटला के सुआगत करत बानी जा।

%%K t07-c2-01-2 p
मई महीना के साँझ के खेत-खरिहान के अँगना केतना मन खुस करे वाला नजारा
लागेला !

%%K t07-c2-02-1 p
२. --- छितिज पर \VUgras{सूरज} \textsuperscript{(1)} धीरे-धीरे डूबत बा, आ
\VUgras{देहात} \textsuperscript{(3)} के आपन लाल \VUgras{किरन}
\textsuperscript{(2)} से भर देत बा। मेहनती \VUgras{हरवाह}
\textsuperscript{(6)} आपन \VUgras{खेत} \textsuperscript{(4)} जोते के जल्दी
करत बा।

%%K t07-c2-02-2 p
आपन \VUgras{हर} \textsuperscript{(5)} से, जवन ताकतवर \VUgras{घोड़ा}
\textsuperscript{(7)} में जोतल बा, ऊ \VUgras{नार} \textsuperscript{(8)}
खींचत बा, जवना में \VUgras{बोवइया} \textsuperscript{(9)} मुट्ठी भर-भर के
\VUgras{बीया} डाली जवना से सोना नियर \VUgras{बाली} उपजी।

%%K t07-c2-03-1 p
३. --- \VUgras{घास काटे वाला} \textsuperscript{(11)}, आपन बड़का हँसुआ
लेके, चरागाह के घास काट लेले बाड़ें, आ सुखावे वाला लोग \VUgras{सूखल
घास} \textsuperscript{(10)} के \VUgras{कँटिया} \textsuperscript{(12)} आ
पंजा से उलट-पलट के सुखा देले बा, आ ई रहल ऊ लोग लवटत बा।

%%K t07-c2-03-2 p
कुछ लोग बैल से खींचल भारी गाड़ी के आगे-आगे चलत बा ; ऊ लोग आपन औजार
कान्ह पर धइले बा। दोसर लोग, जादे अलसी, महकत सूखल घास पर आराम से पसर
गइल बा।

%%K t07-c2-04-1 p
४. --- चलत-चलत ऊ लोग \VUgras{माली} \textsuperscript{(16)} के दोस्ताना
सलाम करत बा, जे \VUgras{फर-बगइचा} \textsuperscript{(13)} में लागल बाड़ें,
\VUgras{फर के गाछ} छाँटत आ \VUgras{नासपाती के गाछ} \textsuperscript{(14)}
में कलम बान्हत। \VUgras{आलूबोखारा के गाछ} \textsuperscript{(15)} में फूल
आवे लागल बा, आ नीक से खाद देहल \VUgras{तरकारी के बगइचा}
\textsuperscript{(17)} में तरकारी, बंदागोभी आ सलाद पहिलही से बढ़िया बा।

%%K t07-c2-04-2 p
छोट भीत पर, सुंदर \VUgras{मोर} \textsuperscript{(18)} आपन पूँछ फइलावत बा,
ताकि लोग ओकर बहुते सुंदर पंख देख के दंग रह जाव।

%%K t07-tit-7 sub
\VUpk{10.2pt}{40}{खेत के अँगना।}
\VUsaut{3.0mm}

%%K t07-c2-05-1 p
५. --- \VUgras{घोड़सार} \textsuperscript{(19)} के केवाड़ खुलल बा आ
\VUgras{पयार} \textsuperscript{(23)} आ चारा के जँगला लउकत बा, जवन ओह
\VUgras{घोड़ी} \textsuperscript{(21)} के ह जेकरा के \VUgras{साईस}
\textsuperscript{(20)} अभिए खोलले बाड़ें। \VUgras{बछेड़ा}
\textsuperscript{(22)}, आपन लमहर गोड़ के साथे एतना मजेदार, कूदत आपन माई
के पाछे-पाछे चलत बा।

%%K t07-c2-06-1 p
६. --- \VUgras{इनार} \textsuperscript{(29)} लगे \VUgras{गाय}
\textsuperscript{(25)} काठ के \VUgras{नाँद} \textsuperscript{(26)} में पानी
पिए जात बाड़ी सँ, जवन ओह लोग के पियासा बुझावे के काम आवेला।
\VUgras{बछड़ा} \textsuperscript{(24)} एह मउका के फायदा उठा के दूध पियत
बा, आ \VUgras{साँड़} \textsuperscript{(27)}, चारा के महक सूँघ के, खुसी से
डकरत बा आ आपन ताकतवर \VUgras{सींग} \textsuperscript{(28)} वाला कपार सान
से उठावत बा।

%%K t07-c2-07-1 p
७. --- सब लोग काम में लागल बा। \VUgras{नौकरानी} \textsuperscript{(32)}
हाथ में इनार के \VUgras{रस्सी} \textsuperscript{(31)} पकड़ले बाड़ी। ऊ
\VUgras{बाल्टी} \textsuperscript{(30)} से पानी खींचत बाड़ी ताकि जनावरन के
पानी पिआवल जाव। \VUgras{खरिहान} \textsuperscript{(33)} लगे एगो मरद पयार
बटोरत बा, आ \VUgras{घर} \textsuperscript{(34)} के केवाड़ के आगे बुढ़िया
\VUgras{कतवइनी} \textsuperscript{(35)} आपन चरखा आ आपन \VUgras{तकुला} से
\VUgras{सन} भा \VUgras{पटुआ} कातत बाड़ी, जइसे पुरनका जमाना में।

%%K t07-tit-8 sub
\VUsaut{3.0mm}
\VUpk{10.2pt}{40}{किसान।}
\VUsaut{3.0mm}

%%K t07-c2-08-1 p
८. --- \VUgras{किसान} \textsuperscript{(36)} ई सब काम के देखरेख करेलें आ
आपन नौकरन के हुकुम देलें। ऊ कहेलें कि \VUgras{हेंगा} \textsuperscript{(40)}
आ \VUgras{कुदाल} \textsuperscript{(39)} के छत के नीचे धर दिहल जाव, जवना
के लोग \VUgras{कुकुर के घर} \textsuperscript{(38)} लगे भुला गइल बा, जवन
\VUgras{कुकुर} \textsuperscript{(37)} के ह।

%%K t07-c2-09-1 p
९. --- ओकर हल्ला से \VUgras{कबूतर} \textsuperscript{(41)} डेरा जात बा आ
पूरा पाँख फइला के \VUgras{कबूतरखाना} \textsuperscript{(42)} में उड़ जात
बा, आ \VUgras{मुरगी} \textsuperscript{(43)} भी, जे जल्दी से
\VUgras{मुरगीखाना} \textsuperscript{(44)} में लवट जात बाड़ी सँ।

%%K t07-c2-10-1 p
१०. --- \VUgras{भेड़सार} \textsuperscript{(48)} के आगे, ई रहलें बुढ़उ
\VUgras{गड़ेरिया} \textsuperscript{(45)} जे आपन \VUgras{रेवड़}
\textsuperscript{(49)} लवटा के ले आवत बाड़ें। आपन \VUgras{लाठी}
\textsuperscript{(46)} थामले, जवन कबो ओकरा से अलग ना होखे, आ आपन रंग
उड़ल \VUgras{कुरता} \textsuperscript{(47)} पहिनले, ऊ बथान की ओर
\VUgras{मेढ़ा} \textsuperscript{(50)}, \VUgras{भेड़} \textsuperscript{(53)},
\VUgras{मेमना} \textsuperscript{(54)}, \VUgras{बकरी} \textsuperscript{(51)}
आ \VUgras{पठरू} \textsuperscript{(52)} हाँकत बाड़ें। ओकर बफादार
\VUgras{कुकुर} \textsuperscript{(55)} आरगुस सबसे पाछे आवत बा आ आपन भूँक
से पिछड़त लोग के जल्दी करावत बा।

%%K t07-c2-11-1 p
११. --- ई सब सोर आ हलचल ओह भल \VUgras{दुधारू गाय} \textsuperscript{(56)}
के सांति ना बिगाड़े, जे आपन \VUgras{घंटी} \textsuperscript{(57)}
बजावत बाड़ी जब ओकरा के \VUgras{गोहाल} \textsuperscript{(58)} के केवाड़ के
आगे दुहल जात बा।

%%K t07-c2-12-1 p
१२. ओकरा जइसे समझ बा कि ओकरा आपन दूध देवे के बा, ताकि
\VUgras{किसानिन} \textsuperscript{(60)} \VUgras{मथनी} \textsuperscript{(61)}
में मक्खन निकाल सकीं, भा ऊ बढ़िया \VUgras{पनीर} \textsuperscript{(64)}
बना सकीं जवन \VUgras{कठौत} \textsuperscript{(63)} पर धइल तखता से चूअत बा।

%%K t07-c2-13-1 p
१३. --- पूरा \VUgras{दूध-घर} \textsuperscript{(59)} के \VUgras{खेत के
नौकर} \textsuperscript{(65)} बहुते साफ रखेलें, जे अभिए ओह बड़का बाल्टी
में \VUgras{दूध के मटकी} \textsuperscript{(62)} धो के आइल बाड़ें जवन ऊ
हाथ में लटकवले बाड़ें।

%%K t07-tit-9 sub
\VUsaut{3.0mm}
\VUpk{10.2pt}{40}{मुरगीखाना।}
\VUsaut{3.0mm}

%%K t07-c2-14-1 p
१४. --- आपन मोटका काठ के जूता पहिनले ऊ थोड़ा भारी चाल से चलेलें, आ
एहसे \VUgras{पीरू} \textsuperscript{(68)}, \VUgras{मादा बतख}
\textsuperscript{(66)}, \VUgras{नर बतख} \textsuperscript{(67)} आ
\VUgras{हंस} \textsuperscript{(69)} भाग जालें, जे आपन \VUgras{चोंच}
\textsuperscript{(70)} चाकर खोल के बेचैन होके चिचिआवेलें।

%%K t07-c2-14-2 p
\VUgras{मुरगा} \textsuperscript{(71)}, जे जादे लड़ाकू बा, आपन लाल
\VUgras{कलगी} \textsuperscript{(72)} तान देला, आपन \VUgras{पूँछ}
\textsuperscript{(73)} के पंख झाड़ेला आ जोर से « कुकड़ूँ-कूँ » बोलेला।

%%K t07-c2-15-1 p
१५. --- \VUgras{मुरगी माई} \textsuperscript{(74)} \VUgras{गोबर के
ढेर} \textsuperscript{(81)} पर आपन \VUgras{चूजा} \textsuperscript{(75)}
के साथे चरत बाड़ी। \VUgras{सूअर के बच्चा} \textsuperscript{(78)} आपन माई ओर
भागत बा, ऊ भारी-भरकम \VUgras{सूअरी} \textsuperscript{(77)} जेकरा के हमनी
के सूअरखाना लगे देखत बानी जा।

%%K t07-c2-16-1 p
१६. --- आपन खाना में \VUgras{खरगोस} \textsuperscript{(79)} ऊ नरम
\VUgras{घास} \textsuperscript{(80)} चाव से खात बाड़ें जवन किसान के बेटी
ओह लोग खातिर ले आवेली। योआनिना के आपन खरगोस बहुते नीक लागेलें आ रोज,
मुरगीखाना से ताजा \VUgras{अंडा} \textsuperscript{(76)} बटोर के, ऊ आपन
छोट दोस्तन से भेंट करे आवेली।
